% !TEX root = ../FinalReport.tex
\chapter{GenAI Safety and Reflection Report}

\textit{Scope and evidence.} GTM Agent supports EU-procurement business development by
ingesting public TED notices and historical awards, ranking candidate suppliers, and
drafting outreach for human review. This report draws on the implemented system, the
ranking backtest, and a five-participant user-acceptance test (UAT), and keeps
implemented controls, evaluation findings, and residual boundaries distinct so that
design intent is not mistaken for evidence of safety.

\section{System boundary and evaluation lens}

\textbf{Where GenAI is used.} GTM Agent does not hand the whole decision chain to one
model. GenAI performs three bounded tasks: extracting soft fields into a schema,
drafting emails from a structured evidence object, and orchestrating a restricted tool
set. Supplier ranking stays outside the generative path (multilingual embeddings,
deterministic rules, and logistic regression), keeping the recommendation measurable
and explainable while confining the model to language and orchestration.

\textbf{What counts as evidence.} The vision specified a two-sided MVP, a hybrid
matcher, feedback and retraining, reviewable drafts, an agent, a ranking backtest, and
UAT. The final model, trained on the earlier 16,812 awards and evaluated on 500
chronologically later held-out awards, reached Hit@1 / Hit@3 / Hit@5 of
44.2\% / 54.0\% / 57.6\%, meeting the Top-1 ($\geq25\%$) and Top-5 ($\geq55\%$) targets.
Two boundaries qualify that number. The test set is restricted to tenders whose winner
is a company the system already knows, so it measures ranking, not coverage of unknown
firms; but the winner is not exempted from the hard filter, and the 26.6\% of cases in
which the filter removes it are counted as misses rather than excluded. Hit@10 and MRR
were not reported, so no claim is made about them.

\section{Safety of GenAI}

Risks arise from generated outputs, agent actions, untrusted inputs, and operation.

\textbf{Factual integrity of emails.} Emails are where generated content reaches users,
so an inaccurate claim about a supplier, a deadline, or prior work is the primary
output risk --- and fluent, confident, unsupported text is a well-documented failure
mode of generative models rather than an occasional defect (Ji et al., 2023). The
system assembles the tender record, supplier profile, scores, match
explanation, and document excerpts, then asks Claude to draft from that evidence alone;
if the model call, JSON parsing, or validation fails, a deterministic template is built
from the same evidence. Every output is stored as a \texttt{draft} and nothing is sent
automatically. These measures limit unsupported content but do not verify every
statement: the current \texttt{EmailValidator} checks subject, body, and length, not
each claim against a source, so human review remains necessary, and a production
version should require every claim to resolve to an identifiable source.

\textbf{Agent authority.} An early version treated \texttt{dry\_run} as a safety
switch, but it blocked database writes without blocking network calls, so a routine
request could start costly TED polling and document scanning. The revised agent makes
dry-run network-free, resolves tenders from the local database by default, and removes
bulk dossier download entirely. Live ingestion is bounded (look-back window, row limit,
once per goal, with an extra condition on award ingestion); the database tool is
read-only (write verbs blocked, row cap, read-only connection); and \texttt{ask\_user}
pauses on ambiguity rather than guessing. The operator console (scheduler, retraining)
is still unauthenticated and would need access control and an audit policy before
deployment.

\textbf{Untrusted inputs, provenance, and privacy.} Tender and dossier text is
attacker-influenced and may carry instructions unrelated to the task --- the indirect
prompt-injection channel described by Abdelnabi et al.\ (2023), in which the payload
reaches the model through retrieved content rather than through the user's own turn.
Schema-constrained extraction, evidence-bounded drafting, restricted tools, and
downstream validation limit --- though do not eliminate --- that channel: untrusted
text cannot alter the deterministic ranking path or gain general write access. The data
is public TED and award records, not private customer data; raw payloads, the evidence
object, and each document's source URL and content hash are retained; files behind
buyer-portal logins are marked \texttt{restricted} rather than scraped. Two
qualifications follow. Supplier contact addresses derived from company websites are
still personal data under Regulation (EU) 2016/679 even though their source is public,
so deployment would need a lawful basis, retention periods, role-based access, and a
privacy review. And because the system generates text that is sent to people on the
system's suggestion, the transparency and human-oversight duties of Regulation (EU)
2024/1689 are the relevant compliance frame for a production version, even though a
lead-generation assistant is not itself a high-risk system under that regulation.

\textbf{Reliability, cost, and fairness.} Rate limiting and retries bound TED
ingestion, local embeddings cut external calls, and a train--serve feature-contract
test guards the deployed matcher, with the deterministic email fallback as an explicit
degradation path. Reliability is not fairness, however: award-history training can
entrench incumbents, and although component scores and adjustable weights make the
basis of a recommendation visible, they do not remove that advantage. With no fairness
metric yet reported, future evaluation should measure ranking quality for smaller and
newer suppliers, not only the recovery of past winners; the ranked group-fairness
criterion of Zehlike et al.\ (2017), which constrains the representation of a protected
group within every prefix of a top-\textit{k} list, is a concrete candidate measure for
that, with firm age or size as the grouping attribute.

\section{Lessons learned and reflection}

\textbf{Delivery versus evidence.} The MVP integrated ingestion, entity resolution,
hybrid matching, evidence-backed drafting, feedback and retraining, the agent, an API,
and a dashboard, but a component existing does not by itself prove product value. The
backtest gives quantitative evidence for the matcher and UAT adds workflow-level
evidence; together they describe completion better than a checklist.

\textbf{What the ranking KPI teaches.} The matcher met its Top-1 and Top-5 targets, yet
raising the training share from 50\% to 80\% did not improve Hit@3 (54.8\% to 54.0\%).
The two runs use different held-out slices, so this is not a controlled comparison and
the roughly one-point difference is within noise; the reading we take from it is
therefore modest --- more data of the same kind did not move the metric. That points to
a representation limit rather than too few examples, since capability profiles are built
from thin award-history text. Enriching those profiles or adding an embedding ensemble
is the more promising direction, and a controlled ablation on a fixed slice would be
needed to confirm it.

\textbf{What the UAT teaches.} Where the backtest asks whether the system retrieves the
known winner, UAT asks whether users can interpret and act on the result. Five
bidder-side participants from transportation, energy, and construction completed
think-aloud sessions covering candidate discovery, weight adjustment, and email
drafting; the full per-participant results are in Part~\ref{part:uat},
Figure~\ref{fig:uat-participants}. Perceived lead quality (3.8/5) exceeded explanation
trust (3.4/5): users saw the business value of the candidates but not always why each
ranked where it did. Showing a score is not the same as explaining it, so a future
interface should tie each recommendation to its CPV, geographic, budget, and semantic
evidence. With only five participants these are design priorities, not population-level
effects.

\textbf{Practices that worked, and what to change.} Assigning ranking to classical ML
and rules kept the path measurable and auditable, and repeated smoke tests against live
TED data caught defects (qualification scoring, budget normalisation) that an
interpretable architecture alone would not. The main organisational lesson is to agree
the interfaces between parallel workstreams early --- the in-app and standalone matchers
were unified late --- and to treat evaluation the same way: each KPI needs a fixed
dataset snapshot, an owner, a reproducible script, and a decision date from the start.
Scope discipline (deferring authenticated scraping, multi-tenancy, blueprint parsing,
and a knowledge graph) kept the MVP focused; the next stage should add claim-to-source
email checks, authenticated operator controls, a larger UAT with a human baseline, and
fairness analysis.

\section{How our understanding of GenAI changed, and conclusion}

The project reframed GenAI as a language-and-orchestration component bounded by evidence
objects, schemas, permissions, and user confirmation --- not a general decision-maker;
the dry-run incident showed that a control matters only when it blocks the exact action
that produces the risk. UAT separated usefulness from trust: a system can surface
valuable leads while users stay unsure of the reasoning, and a fluent draft can still
need factual review, so evaluating a GenAI application means weighing task completion,
ranking quality, factual support, user understanding, operational limits, and
auditability together. GTM Agent now forms an integrated MVP whose safety case rests on
evidence-constrained generation, human review, restricted tools, offline-by-default
agent behaviour, and traceable public data; the next stage should keep testing those
boundaries, especially email-claim verification, operator access control, and broader
user and fairness evaluation.
